% =====================================================================
%  9 -- SILICON BURNING ON REAL HISTORIES, AND THE LABEL PATHOLOGY
%       (source node S13, Part III).
%
%  Source: tier3.md lines 3457-3839 in full and in order --
%          Part III header (3457-3461),
%          III.1 What a real trajectory looks like (3463-3499),
%          III.2 The census (3501-3540),
%          III.3 The mechanism: gh-575 (3542-3617),
%          III.4 The witness experiment (3619-3660),
%          III.4b How this could have been caught on day one (3662-3711),
%          III.5 The consequence: Ye itself is wrong (3713-3743),
%          III.6 The epistemics (3745-3812),
%          III.7 What the pathology does *not* invalidate (3814-3838),
%          III.8 Self-check for S13 (3830-3839; 7 questions).
%
%  Rendering notes:
%    * the Horn A / Horn B blockquote (3764-3774) -> aside.
%    * the two "!" (warning-glyph) list items in III.6 keep the glyph
%      as \warnglyph{}.
%    * the "corrected here in the 2026-08-18 audit" parenthetical in
%      III.6 is kept verbatim.
%  Cross-refs reworded: none (anchor rewrites only; no direction
%  changed -- every referenced section keeps its relative order).
% =====================================================================

\section{Silicon burning on real histories, and the label pathology}
\label{part:pathology}

\textbf{The payoff node, and the one that changed the project.} Everything
before this was instrumentation. Here the instruments get pointed at the
dataset the emulator is supposed to learn, and the dataset fails.

% =====================================================================
\subsection{What a real trajectory looks like}
\label{sec:realtrajectory}

Take one constant-$(T, \rho)$ hydrostatic run at, say, $\Tn = 4$,
$\rho = 10^{8}$, starting from a silicon-rich mix. Read the output rows in
order and you see four phases:

\begingroup\small
\begin{tabularx}{\textwidth}{@{}L{3.6cm} L{2.6cm} Y@{}}
\toprule
\textbf{phase} & \textbf{age} & \textbf{what happens} \\
\midrule
\textbf{1 --- free-particle relaxation} & $10^{-8}$ \ldots\ ${\sim}10^{-6}$ s &
  whatever free n, p, $\alpha$ the initial composition carried are captured
  almost instantly: $\lvert \Delta X\rvert / X \sim 1$ in a single step. A
  \textbf{stiff relaxation}, not a linear step \\
\textbf{2 --- strong-sector equilibration} & ${\sim}10^{-6}$ \ldots\ $10^{0}$ s &
  the fast capture/photodisintegration pairs balance within each group;
  $\kappa$ on those columns falls; the composition organises into the Si and
  Fe groups \\
\textbf{3 --- QSE burning} & $10^{0}$ \ldots\ $10^{4}$--$10^{6}$ s &
  net flow through the bridges carries material from the Si group to the Fe
  group. \textbf{This is the physics the project is about}; \enuc{} is roughly
  steady \\
\textbf{4 --- weak drift} & after arrival &
  strong sector frozen at its attractor; \Ye{} continues to fall through
  electron captures; composition changes slowly \\
\bottomrule
\end{tabularx}
\endgroup

Phase 1 is the single most under-appreciated fact about the \emph{training}
labels. The Sobol initial compositions carry free nucleons (median
$X_{\mathrm{neut}} \approx 1.7 \times 10^{-2}$), so even the \textbf{shortest
label timestep} is not a linear step. The measurement: \textbf{99.99\% /
100.00\% of (state, isotope) cells have $\tau < \mathrm{dt}$} at
$\mathrm{dt}_1 \approx 1.01 \times 10^{-6}$ s, with label $|\Delta X|/X$ median
$\approx 1$, and only \textbf{20 / 0} linearisable cells out of 2.4M / 4.5M
\resultsref{2026-07-10}.

Read that again. \textbf{Every label in the shipped training set encodes a full
stiff relaxation, at every timestep, everywhere in the box.} There is no
``small step'' regime in this dataset. Consequences:

\begin{itemize}
\item Target A's $\Phi$ is a \emph{time-integrated effective flux}
  $\Phi_j = \int \varphi_j\,\dd t$ over the label interval, never an
  instantaneous rate. \code{CLAUDE.md} states this and
  \code{fluxes/integrate.py} is the reference producer.
\item Conditioning the model on $\log \Delta t$ is conditioning on the wrong
  variable; the natural one is $\Delta t/\tau$, since what varies across the
  dataset is \emph{how many relaxation times} the step covers
  (\foreignreg{2}{R-10}).
\item Any handshake between the engine's instantaneous rates and the labels can
  only be attempted on the small subset where $\tau > 10\Delta t$; that is why
  the trajectory handshake is restricted to linearisable cells (median
  rate-level residual $3.6 \times 10^{-3}$ / $2.8 \times 10^{-3}$, $\le 5\%$
  for 90.3\% / 94.4\% of cells \resultsref{2026-07-10}).
\end{itemize}

% =====================================================================
\subsection{The census: how far are the labels from NSE?}
\label{sec:census}

The question that started the investigation is simple. At high temperature,
silicon burning should approach NSE (\S\ref{sec:staircase},
\S\ref{sec:endpoint}): the fast pairs balance, and with enough time the
composition should be the Saha solution $X(T, \rho, \Ye)$. So take the
\textbf{longest} label (dt $= 10^{2}$ s), at \textbf{high} \Tn{}, and compare
against the project's independent Saha solver.

The solver is independent in the sense that matters: it is built from the same
pynucastro nuclear inputs (\citealt{smithclark2023}; \citealt{rauscher1997} /
\citealt{rauscher2003} partition functions, binding energies, spins) but solves
the equilibrium conditions itself, in log space with a damped Newton and a
bisection fallback, and it agrees with pynucastro's own NSE solver to
$\mathbf{2.7 \times 10^{-10}}$ / $\mathbf{1.5 \times 10^{-10}}$ \textbf{dex}
over a 27-state grid \resultsref{2026-07-10}. It is a trustworthy reference.

The result, over \Tn{} bins from $[5.0, 5.5)$ to $[7.0, 7.94)$, 120 states per
bin, \resultsref{2026-07-11}:

\begin{tabularx}{\textwidth}{@{}Y L{4.0cm} L{3.4cm}@{}}
\toprule
\textbf{quantity} & \textbf{\msn{80}} & \textbf{\msn{151}} \\
\midrule
median $\max|\Delta\log_{10}X|$ over $X > 10^{-6}$ species &
  \textbf{6.9 $\to$ 11.5 dex} across bins & \textbf{6.7 $\to$ 13.0 dex} \\
fraction of states more than 1 dex from NSE &
  \textbf{1.000 in every bin} & \textbf{1.000} \\
\bottomrule
\end{tabularx}

\textbf{Not one sampled high-\Tn{} label endpoint is anywhere near NSE.} And
the discrepancy \emph{grows} with temperature, which is the opposite of the
physical expectation --- hotter should mean faster equilibration.

The displaced states have structure, which is the first clue that this is a
mechanism rather than noise:

\begin{itemize}
\item \textbf{high $\rho$}: \nuc{Si}{30} / \nuc{Si}{29} / \nuc{Mg}{26}
  dominated. State 80 ($\Tn = 7.08$, $\rho = 3.0 \times 10^{8}$): si30
  $X = 0.578$, where NSE says fe56 $= 0.72$.
\item \textbf{low $\rho$}: \nuc{C}{12} / \nuc{O}{16} dominated. State 646
  ($\Tn = 7.78$, $\rho = 2.0 \times 10^{7}$): c12 $= 0.294$ + o16 $= 0.266$,
  where NSE says \nuc{He}{4} and the Fe group.
\end{itemize}

And these are \textbf{the same class of state} as the Step-5 trajectory-stall
frozen compositions (\nuc{Si}{30}-dominated at $\Tn = 6.1$). Two independent
investigations, from different directions, landed on the same anomaly.

% =====================================================================
\subsection{The mechanism: gh-575, traced from a Fortran branch}
\label{sec:gh575}

The cause is a bug in MESA r23.05.1's detailed-balance reverse-rate
construction --- MESAHub/mesa issue \textbf{gh-575} \citep{gh575} (opened 2023-08-07, closed
2024-05-06; the fix, PR \#632, merged May 2024 and first shipped in release
24.08.1 --- verified by measurement, \resultsref{2026-07-10}; the release-note
text is recorded in \resultsref{2026-07-09}: ``incorrect phase space factors
for reverse reaction rates involving greater than 2 reactants or products
\ldots\ inconsistent equilibrium compositions \ldots\ at temperatures exceeding
4 GK'').

\paragraph{The static locus.} In \code{rates/private/reaclib\_support.f90},
\code{compute\_rev\_ratio} applies the detailed-balance phase-space factor
%
\begin{equation}\label{eq:revfac}
\mathrm{fac} \;=\; \frac{1}{N_A}\left(\frac{10^{9}k_B}{2\pi\hbar^{2}N_A}\right)^{3/2} \;=\; 9.868\times10^{9},
\qquad \log_{10}\mathrm{fac} = 9.994,
\end{equation}
%
together with its accompanying $T^{3/2}$, \textbf{only in the single-product
branch} (\code{reaclib\_support.f90:197, 222--235}; the constant is
\citealt{rauscher2000}'s $9.8685 \times 10^{9}\,\Tn^{3/2}$, recomputed
$9.86846 \times 10^{9}$ with CODATA 2018, \citealt{tiesinga2021}). Tier~1
\S III derives why that factor is there: the reverse of an
A + B $\to$ C reaction involves a change $\Delta N$ in the number of free
particles, and each unit of $\Delta N$ brings a factor of
$(2\pi\mu kT/h^{2})^{3/2}/(\rho N_A)$ from the phase-space integration. A
reverse rate with more than one product needs the factor raised to the power
$|\Delta N|$; the buggy code applies it once, or not at all.

\paragraph{The empirical size.} Measured against pynucastro at
$\Tn \in \{1.6, 4.0, 7.9\}$ \resultsref{2026-07-09}:

\begin{tabularx}{\textwidth}{@{}Y L{5.0cm}@{}}
\toprule
\textbf{channel class} & \textbf{$|\Delta\log_{10}|$} \\
\midrule
$|\Delta N| = 1$ multi-body inverses & \textbf{10.0 -- 11.1 dex} \\
$|\Delta N| = 2$ (h1+h1+he4+he4 $\to$ he3+be7) & \textbf{20.6 -- 22.7 dex} \\
chapter-8 photodisintegration reverses (1 $\to$ 3) &
  \textbf{low by 9.5 -- 11.3 dex} \\
harness control: all matched clean REACLIB forwards &
  median \textbf{0.0 exactly} \\
\bottomrule
\end{tabularx}

The last row is what makes the rest credible: the comparison harness
reproduces clean channels bit-for-bit, so the 10-dex discrepancies are not the
harness.

\textbf{The chapter-8 row is the one that matters here, and it goes beyond what
the NNN paper's own Appendix B names.} The paper names one light-sector channel
(n + n + \nuc{He}{4} + \nuc{He}{4} $\to$ \nuc{H}{3} + \nuc{Li}{7}) and states
the class --- ``all reactions involving more than two reactants and/or
products'' --- which contains the 1 $\to$ 3 reverses without identifying them,
and it states that \msn{80}'s equilibrium composition is ``qualitatively
consistent'' with the larger networks; this project found that the 1 $\to$ 3
photodisintegration reverses --- including \textbf{c12 $\to$ 3$\alpha$, the
reverse of the triple-$\alpha$ reaction}, and be9 $\to$ n + 2$\alpha$, li6
$\to$ n + p + $\alpha$ --- are low by 9.5--11.3 dex in stock r23.05.1. Those
reactions control \textbf{light $\leftrightarrow$ heavy equilibration}: they
are how the free-$\alpha$ reservoir talks to the carbon/oxygen sector.

Now the mechanism is visible. If c12 $\to$ 3$\alpha$ is suppressed by ten
orders of magnitude while 3$\alpha$ $\to$ c12 is correct, then detailed balance
is broken by ten orders of magnitude on the reaction that links the $\alpha$
reservoir to everything above it. The system relaxes to a \textbf{fixed point
of the broken rate set} --- an equilibrium of a rate network that satisfies no
thermodynamic constraint. At high $\rho$ that fixed point is
\nuc{Si}{30}-dominated; at low $\rho$ it is \nuc{C}{12}/\nuc{O}{16}-dominated,
because carbon cannot be photodisintegrated back into $\alpha$ particles.

\textbf{This also explains the $\kappa$ floor.} Tier~2 \S IV.8b's Wegscheider
measurement (the cycle condition of \citealt{wegscheider1901}, in the
generalised form of \citealt{schuster1989}) found
$Q \notin \mathrm{rowspace}(\nu)$ for the strong sector and predicted a
$\kappa$ data floor of $6 \times 10^{-4}$ \ldots\ $4 \times 10^{-3}$; the
measured stock-MESA $\kappa$ floor at NSE is median $3.6 \times 10^{-3}$ /
$3.3 \times 10^{-3}$ \resultsref{2026-07-10}, and \textbf{$\kappa \to 1.0$
exactly on gh-575 channels} --- a broken pair is maximally out of balance by
construction.

% =====================================================================
\subsection{The witness experiment}
\label{sec:witness}

A mechanism is a story until it makes a prediction that could fail. The
prediction: \emph{if gh-575 is the cause, then running the exact same label
states under the exact same configuration but with the fixed MESA should relax
toward NSE instead.}

\code{scripts/step6\_label\_nse\_census.py} ran it on the two witness states,
with bbq built against both MESA versions, under exact label conditions and
initial compositions \resultsref{2026-07-11}:

\begingroup\small
\begin{tabularx}{\textwidth}{@{}L{4.6cm} Y Y@{}}
\toprule
 & \textbf{state 80 (\Tn{} 7.08, $\rho$ 3.0e8)} &
   \textbf{state 646 (\Tn{} 7.78, $\rho$ 2.0e7)} \\
\midrule
\textbf{shipped label} & si30 $= 0.5615$ & c12 0.294 + o16 0.266 \\
\textbf{stock r23.05.1 bbq} & reproduces the label to \textbf{0.01 dex} &
  reproduces to \textbf{0.03 dex} \\
\textbf{MESA 24.08.1 bbq} (gh-575 verified fixed) &
  \textbf{fe56 $= 0.985$}, 1.2 dex residual, \Ye{} still evolving &
  \textbf{he4 $= 0.521$ + fe56 $= 0.378$} \\
\textbf{independent pf-true reference integrator} &
  fe56 $= 0.70$ at $10^{-6}$ s (vs 24.08.1's 0.66) & --- \\
\bottomrule
\end{tabularx}
\endgroup

Three independent codes --- stock bbq, fixed bbq, and this project's own BDF
integrator with pf-corrected reverse rates --- arrange themselves exactly as
the mechanism predicts. Stock reproduces the pathology; fixed and independent
both relax toward NSE.

\textbf{This is the strongest single piece of evidence in the whole of Phase
0}, and it is worth naming why. It is a \emph{controlled} experiment: same
inlist, same initial composition, same solver tolerances, same screening, same
weak tables --- one code version changed (whose relevant diff is the gh-575
fix; 24.08.1 also carries a minor spillover into approx21), and the outcome
switches from ``reproduces the shipped labels'' to ``does the physically
expected thing''. Everything else in this project is a comparison of computed
numbers; this is an intervention.

\textbf{The dt-freeze signature} seals it. State 80's final composition is
\emph{identical to four digits} (si30 $= 5.615 \times 10^{-1}$) across
\textbf{dt $= 10^{-5}$ \ldots\ $10^{1}$ s}, and is already reached at
dt $= 10^{-6}$ \resultsref{2026-07-11}. A slow transient would show dt
dependence. A fixed point does not. The labels are sitting at an attractor of
the label generator, reached almost immediately and held for seven decades of
time.

% =====================================================================
\subsection{How this could have been caught on day one}
\label{sec:dayone}

Worth extracting as a reusable test, because the pathology was found in Step 6
and could have been found in Step 1 with an hour's work and no MESA at all.

\textbf{The generic test: is the label set a fixed point of its own generator,
and is that fixed point the one physics predicts?}

For a sample of high-temperature states:

\begin{enumerate}
\item take the label at the \textbf{longest} $\Delta t$;
\item compare it against an independent equilibrium reference --- the Saha
  solution $X_{\mathrm{NSE}}(T, \rho, \Ye)$;
\item compare the label \textbf{across $\Delta t$ decades} --- the freeze
  signature.
\end{enumerate}

Step 2 needs a Saha solver, which pynucastro ships. Step 3 needs nothing but
the shipped CSVs --- and it is the \emph{decisive} one: a composition identical
to four digits across dt $= 10^{-5}$ \ldots\ $10^{1}$ s is a fixed point, full
stop. \textbf{That check reads nine files and takes a minute.} It would have
fired on day one.

Three general lessons, stated so they transfer:

\paragraph{(i) A benchmark dataset is a fixed point of its generator, and that
is a testable claim about the generator, not about nature.} Any dataset
produced by running a solver to completion encodes the solver's attractors.
Checking those attractors against an independent equilibrium theory is the
cheapest possible audit of a labelled dataset, and it is not, to our knowledge,
standard practice in scientific ML.

\paragraph{(ii) Invariance across a control parameter is a stronger signal than
agreement with a reference.} The dt-freeze is more convincing than the NSE
distance, because it needs no reference at all --- it is internal to the data.
Look for quantities that \emph{should} vary and do not.

\paragraph{(iii) The check that fires is the one comparing to something the
pipeline does not share.} Tier~1 \S VIII.D's habit --- \emph{ask of any green
test what would still be wrong if it passed} --- generalises to \emph{any check
whose two sides share a provenance can only find typos.} The rate cross-check
(Tier~1 \S VI) compared MESA against pynucastro and found the bug's
\emph{static} form in Step 4; it took an independent \textbf{equilibrium}
reference, which shares no code with either, to find the bug's \emph{dynamical
consequence}.

Registered as \reg{T-29}: adopt the fixed-point audit as a standing check on
any labelled dataset the project ingests, including its own future $\Phi$
labels.

% =====================================================================
\subsection{The consequence: \texorpdfstring{\Ye}{Ye} itself is wrong}
\label{sec:yewrong}

If the bug only displaced \emph{which} silicon isotope dominates, it would be a
composition-detail problem. It is not, and this is the finding that forced an
escalation.

At state 80, over $10^{2}$ s \resultsref{2026-07-11}:

\begin{tabularx}{\textwidth}{@{}L{5.0cm} Y@{}}
\toprule
\textbf{path} & \textbf{\Ye{} evolution} \\
\midrule
labels / stock r23.05.1 & 0.4717 $\to$ \textbf{0.4713}
  ($\Delta = -4 \times 10^{-4}$) \\
fixed MESA (24.08.1) & 0.4717 $\to$ \textbf{0.4619}
  ($\Delta = -9.8 \times 10^{-3}$) \\
\bottomrule
\end{tabularx}

\textbf{A factor of ${\sim}25$ in the \Ye{} change} (24.5 at the quoted
4-digit precision). The reason is direct: electron-capture rates are
per-nucleus and species-specific (\S\ref{sec:yecarriers}; \citealt{heger2001}
tabulate the dominant EC flows; \citealt{langanke2000} give the iron-peak
rates), so \emph{which species host the material} determines which EC channels
are active. A composition frozen at si30-dominated hosts EC on silicon-group
nuclei; a composition relaxed to fe56-dominated hosts EC on the iron peak,
where the rates are much larger and the thresholds much more favourable. The
bug displaces the composition, the composition selects the EC channels, and the
EC channels set $\dd\Ye/\dd t$.

Put in the terms of \S\ref{sec:gatemap}'s budget: the \textbf{discrepancy
between the label \Ye{} evolution and the physically-correct one is
${\sim}9 \times 10^{-3}$ over $10^{2}$ s} --- which is larger than the
project's entire per-trajectory \Ye{} budget ($5 \times 10^{-3}$) and three
thousand times the per-step gate. At $\Tn \gtrsim 5$, the labels are not merely
imprecise; they are outside the tolerance the emulator is being asked to
achieve, in the quantity the emulator exists to predict.

% =====================================================================
\subsection{The epistemics: what a negative result about your own dataset obliges}
\label{sec:epistemics}

This is the part of Tier 3 that is about method rather than physics, and it is
the part most worth internalising.

\paragraph{What was superseded.} Step 4 had established a working premise: the
training labels were assumed to have been generated with the authors'
\emph{locally fixed} MESA --- an inference from the paper's Appendix B, which
reports ``We fixed the bug in the current mesa version'' without stating which
build produced the training sets, and whose software list names mesa-23.05.1
--- so labels are clean and \emph{our} stock MESA is the outlier. That premise
was reasonable, an inference from the paper rather than a statement in it, and
recorded (RESULTS 2026-07-09 phrases it as ``the paper's App. B says so'',
which over-reads the appendix in the same way; corrected here in the 2026-08-18
audit). It is now known to be false at the fixed-point level for
$\Tn \gtrsim 5$ --- because the chapter-8 1 $\to$ 3 reverses this project found
``beyond the paper'' are demonstrably label-active, whatever the paper's own
channel list says. The \code{RESULTS.md} entry states the supersession
explicitly and identifies which earlier row it overrides. \textbf{That is the
discipline: a superseding measurement names its predecessor.}

\paragraph{What was validated by accident.} \code{ADR 0005} chose stock
r23.05.1 for the rerun campaign on \emph{comparability} grounds, before the
pathology was known. The pathology measurement then showed stock reproduces
label behaviour to 0.01 dex --- so the choice was measured-correct after the
fact. A decision made for one good reason turning out to be right for a
stronger one is luck, and it should be recorded as luck rather than as
foresight.

\paragraph{What it forces.} Five consequences were recorded
\resultsref{2026-07-11}, and two of them are \emph{escalations to a human}, not
technical fixes:

\begin{enumerate}
\item rerun-campaign comparability \textbf{validated} (above);
\item the Step-5 stall anomaly \textbf{reinterpreted} --- the frozen states are
  the bug's displaced pseudo-equilibria; the output-dt hypothesis is at most
  secondary (\S\ref{sec:stall});
\item \code{validate\_campaign}'s terminal-NSE check must \textbf{expect} the
  displaced equilibrium on stock reruns at $\Tn \gtrsim 5$, with 24.08.1 as
  the physics-true witness (\S\ref{sec:validation}, check 2);
\item \warnglyph{}~\textbf{Phase-1 supervision at $\Tn \gtrsim 5$ faces a
  benchmark-versus-physics fork} --- escalated, human decision;
\item \warnglyph{}~\textbf{NNN itself was trained on these labels} --- an
  external-communication decision, escalated.
\end{enumerate}

\textbf{The fork (item 4) is a genuine dilemma, and it is worth stating both
horns fairly.}

\begin{aside}
\textbf{Horn A --- train on the labels.} The emulator's job is to reproduce
the reference solver. The published benchmark (NNN) was trained on these
labels, and comparability requires matching them. Training on physics-true
labels at $\Tn \gtrsim 5$ would make every published comparison in that band
meaningless --- the emulator would ``lose'' against a baseline that is
reproducing a bug.

\textbf{Horn B --- train on corrected physics.} The emulator's job is to be a
fast nuclear network. A network that reproduces a known bug is a liability the
moment anyone uses its output for science, and the bug is \emph{in the quantity
the project exists to predict} (\S\ref{sec:yewrong}). Every downstream \Ye{}
conclusion in that band would be wrong by design.
\end{aside}

There is a third option that the study notes should name because it is
probably right and nobody has written it down: \textbf{declare a validity
domain.} $\Tn < 5$ is where local $\Phi$ labels are proven feasible
(\S\ref{sec:phisupervision}), where every kill-test gate passes, and where the
labels are not pathological. Restricting the emulator's claimed domain to
$\Tn < 5$ and handing $\Tn \ge 5$ to an NSE table or the real solver (which is
the \S\ref{sec:landscape} hybrid anyway, and exactly what an OOD-fallback gate
is \emph{for}) resolves the fork by refusing it. It costs a claim about the hot
band and buys a defensible product. Registered as \reg{T-03}.

\textbf{Item 5 is the uncomfortable one.} The published NNN model, its reported
accuracies, and this project's own reproduction of them
(\resultsref{2026-07-08}, ratio 1.000 in all 108 comparisons) are all measured
against labels that are wrong above 5 GK. That does not make the NNN paper's
\emph{methods} wrong --- its errors are measured against its own labels, which
is the standard and correct protocol --- but it does mean that a statement like
``the NNN reproduces the composition to X\%'' carries an unstated qualifier in
that band. How and whether to communicate that externally is a judgement call
for a person, and the project correctly escalated it rather than deciding it
inside a results log.

% =====================================================================
\subsection{What the pathology does \emph{not} invalidate}
\label{sec:notinvalidated}

Being precise about the blast radius matters, because a finding this large
invites over-generalisation.

\textbf{Still valid:}

\begin{itemize}
\item \textbf{Everything below $\Tn \approx 5$.} The gh-575 channels are
  light-sector and chapter-8 multi-body; they matter where the
  $\alpha$/nucleon reservoir controls equilibration, which is the hot regime.
  The QSE window (3.3--5 GK) shows integrator/label agreement of 0.74--0.99 in
  band \resultsref{2026-07-12}.
\item \textbf{The kill-test itself.} It is a measurement \emph{on the label
  manifold}, and the reruns are exactly label-comparable (0.996/0.998
  early-time agreement). The question ``is Target A viable on this data'' is
  well-posed regardless of whether the data is physically right --- because
  that is the data the model will see. \S\ref{sec:verdict} develops this.
\item \textbf{All the algebra.} Conservation, the projector, the flux engine's
  rate agreement, the Saha solver --- none of it depends on the labels being
  right.
\item \textbf{The bridge sets and the \nuc{Sc}{45} confirmation}, which live
  in the QSE window.
\end{itemize}

\textbf{Compromised:}

\begin{itemize}
\item Any \emph{physical} claim about equilibrium structure at $\Tn \gtrsim 5$
  on this data.
\item The Guidry $\varepsilon$-sweep result read as a statement about silicon
  burning (as opposed to about the label manifold) --- see \S\ref{sec:gate5},
  which is careful about exactly this.
\item \Ye{} supervision in the hot band (\S\ref{sec:yewrong}).
\end{itemize}

% =====================================================================
\subsection{Self-check}
\label{sec:selfcheck-pathology}

\begin{selfcheck}
% source III.8 Q1
\item Why does the \emph{shortest} label timestep in the shipped set encode a
  full stiff relaxation? Give the measurement.
% source III.8 Q2
\item Reproduce the argument from ``a phase-space factor applied only in the
  single-product branch'' to ``the labels sit at a
  \nuc{C}{12}/\nuc{O}{16}-dominated fixed point at low density''.
% source III.8 Q3
\item The census found the NSE distance \emph{growing} with temperature. Why is
  that the opposite of the physical expectation, and why is it the signature
  of a mechanism rather than of noise?
% source III.8 Q4
\item Describe the witness experiment, name its one changed variable, and say
  what result would have falsified the mechanism.
% source III.8 Q5
\item The bug displaces composition. Explain in two steps why it therefore
  corrupts \Ye{}, and give the measured factor.
% source III.8 Q6
\item State both horns of the benchmark-versus-physics fork fairly, then give
  the third option and what it costs.
% source III.8 Q7
\item Name three things the pathology does not invalidate, with the reason for
  each.
\end{selfcheck}
